% --- Archon physics formalization source begin ---
% archon:physics
% archon:covers PhyXMiniProblems/problem_phyx_mini_0238.lean
% archon:source-report reports/phyx_mini/problem_phyx_mini_0238.source.json

\chapter{Physics problem phyx\_mini\_0238}
\label{ch:PhyXMiniProblems_problem_phyx_mini_0238}

\paragraph{Problem source.}
\#\# Physical scenario

A light, cubical container of volume \textbackslash{}( a\textasciicircum{}3 \textbackslash{}) is initially filled with a liquid of mass density \textbackslash{}( \textbackslash{}rho \textbackslash{}) as shown in figure. The cube is initially supported by a light string to form a simple pendulum of length \textbackslash{}( L\_i \textbackslash{}), measured from the center of mass of the filled container, where \textbackslash{}( L\_i \textbackslash{}gg a \textbackslash{}). The liquid is allowed to flow from the bottom of the container at a constant rate \textbackslash{}( dM/dt \textbackslash{}). At any time \textbackslash{}( t \textbackslash{}), the level of the liquid in the container is \textbackslash{}( h \textbackslash{}) and the length of the pendulum is \textbackslash{}( L \textbackslash{}) (measured relative to the instantaneous center of mass) as shown in figure.The liquid completely runs out of the container.

\#\# Question

What is the period of the pendulum?

\#\# Answer choices

- A: A: \$2\textbackslash{}pi \textbackslash{}sqrt\{\textbackslash{}frac\{L\_i\}\{g\textasciicircum{}2\}\}\$

- B: B: \$\textbackslash{}pi \textbackslash{}sqrt\{\textbackslash{}frac\{L\_i\}\{g\}\}\$

- C: C: \$2\textbackslash{}pi \textbackslash{}sqrt\{\textbackslash{}frac\{g\}\{L\_i\}\}\$

- D: D: \$2\textbackslash{}pi \textbackslash{}sqrt\{\textbackslash{}frac\{L\_i\}\{g\}\}\$

\#\# Figure caption (auxiliary; use the image as primary evidence)

The image consists of two diagrams showing a block suspended by a rope. 

In the left diagram:
- A cube-shaped block is hanging from a horizontal surface by a single rope.
- The block has sides labeled as "a".
- The length of the rope from the surface to the block is marked as "Li".

In the right diagram:
- The same block is now positioned with a partial distance "h" below the starting point shown on the left.
- The rope length from the surface to the block is marked as "L".
- There is a horizontal line marking the distance "h" which shows the displacement from the left position.

Both diagrams illustrate the concept of a hanging block with specified dimensions and hanging lengths.

\#\# Dataset metadata

Resonance and Harmonics; Physical Model Grounding Reasoning, Multi-Formula Reasoning

\paragraph{Recorded answer/context.}
D: D: \$2\textbackslash{}pi \textbackslash{}sqrt\{\textbackslash{}frac\{L\_i\}\{g\}\}\$

\paragraph{Figure/image path.}
/root/proposal\_for\_physic/science-mango/phyx\_mini\_run/phyx\_data/test\_image/238.png

\paragraph{Formalization target.}
create a compiling Lean file with sorry bodies at `PhyXMiniProblems/problem\_phyx\_mini\_0238.lean`. The Lean declarations must preserve the physical quantities, dimensions or dimensional roles, figure labels, governing-law hypotheses, and final relation expressed by this problem.
Use Mathlib/Physlib names found through LeanExplore where available. If a physics API is missing, introduce faithful local abstractions rather than scalar placeholder aliases.

\begin{theorem}[Physics formalization target]
\label{thm:physics:phyx_mini_0238:target}
\lean{PhyXMiniProblems.ProblemPhyXMini0238.period_of_leaking_cubical_pendulum}
\uses{def:physics:phyx-mini-0238:phyxminiproblems-problemphyxmini0238-lengthinmeters, def:physics:phyx-mini-0238:phyxminiproblems-problemphyxmini0238-accelerationinmeterspersecondsquared, def:physics:phyx-mini-0238:phyxminiproblems-problemphyxmini0238-timeinseconds, def:physics:phyx-mini-0238:phyxminiproblems-problemphyxmini0238-leakingcubependulumsetup, def:physics:phyx-mini-0238:phyxminiproblems-problemphyxmini0238-matchesprimaryfigure, def:physics:phyx-mini-0238:phyxminiproblems-problemphyxmini0238-matchesproblemdata, def:physics:phyx-mini-0238:phyxminiproblems-problemphyxmini0238-duringemptying, def:physics:phyx-mini-0238:phyxminiproblems-problemphyxmini0238-hasphysicalparameters, def:physics:phyx-mini-0238:phyxminiproblems-problemphyxmini0238-linearizedsimplependulumperiodseconds, def:physics:phyx-mini-0238:phyxminiproblems-problemphyxmini0238-obeysleakingsimplependulumlaws, def:physics:phyx-mini-0238:phyxminiproblems-problemphyxmini0238-answerchoice, def:physics:phyx-mini-0238:phyxminiproblems-problemphyxmini0238-answerchoice-displayedscalar, def:physics:phyx-mini-0238:phyxminiproblems-problemphyxmini0238-recordedanswerchoice, def:physics:phyx-mini-0238:phyxminiproblems-problemphyxmini0238-smallangleperioderrorbudgetseconds, def:physics:phyx-mini-0238:phyxminiproblems-problemphyxmini0238-longstringperioderrorbudgetseconds}
The assigned autoformalize agent should translate this physics problem into Lean declarations in the covered file, with theorem and lemma proof bodies written as `by sorry`.
\end{theorem}
\begin{proof}
This is an autoformalization task, not a proof task. Produce statements that can later be proved by the physics prover without weakening the physical model.
\end{proof}

% --- Archon declaration topology begin ---
\section{Declaration topology}
\begin{definition}[volume Dimension]
  \label{def:physics:phyx-mini-0238:phyxminiproblems-problemphyxmini0238-volumedimension}
  \lean{PhyXMiniProblems.ProblemPhyXMini0238.volumeDimension}
  The physical dimension of volume, L³
\end{definition}
\begin{proof}
  This is the declared model definition of volume Dimension.
\end{proof}
\begin{definition}[mass Density Dimension]
  \label{def:physics:phyx-mini-0238:phyxminiproblems-problemphyxmini0238-massdensitydimension}
  \lean{PhyXMiniProblems.ProblemPhyXMini0238.massDensityDimension}
  The physical dimension of mass density, M L⁻³
\end{definition}
\begin{proof}
  This is the declared model definition of mass Density Dimension.
\end{proof}
\begin{definition}[mass Flow Rate Dimension]
  \label{def:physics:phyx-mini-0238:phyxminiproblems-problemphyxmini0238-massflowratedimension}
  \lean{PhyXMiniProblems.ProblemPhyXMini0238.massFlowRateDimension}
  The physical dimension of mass flow rate, M T⁻¹
\end{definition}
\begin{proof}
  This is the declared model definition of mass Flow Rate Dimension.
\end{proof}
\begin{definition}[acceleration Dimension]
  \label{def:physics:phyx-mini-0238:phyxminiproblems-problemphyxmini0238-accelerationdimension}
  \lean{PhyXMiniProblems.ProblemPhyXMini0238.accelerationDimension}
  The physical dimension of acceleration, L T⁻²
\end{definition}
\begin{proof}
  This is the declared model definition of acceleration Dimension.
\end{proof}
\begin{definition}[Length Quantity]
  \label{def:physics:phyx-mini-0238:phyxminiproblems-problemphyxmini0238-lengthquantity}
  \lean{PhyXMiniProblems.ProblemPhyXMini0238.LengthQuantity}
  A nonnegative physical length
\end{definition}
\begin{proof}
  This is the declared model definition of Length Quantity.
\end{proof}
\begin{definition}[Volume Quantity]
  \label{def:physics:phyx-mini-0238:phyxminiproblems-problemphyxmini0238-volumequantity}
  \lean{PhyXMiniProblems.ProblemPhyXMini0238.VolumeQuantity}
  \uses{def:physics:phyx-mini-0238:phyxminiproblems-problemphyxmini0238-volumedimension}
  A nonnegative physical volume
\end{definition}
\begin{proof}
  This is the declared model definition of Volume Quantity.
\end{proof}
\begin{definition}[Mass Quantity]
  \label{def:physics:phyx-mini-0238:phyxminiproblems-problemphyxmini0238-massquantity}
  \lean{PhyXMiniProblems.ProblemPhyXMini0238.MassQuantity}
  A nonnegative physical mass
\end{definition}
\begin{proof}
  This is the declared model definition of Mass Quantity.
\end{proof}
\begin{definition}[Mass Density Quantity]
  \label{def:physics:phyx-mini-0238:phyxminiproblems-problemphyxmini0238-massdensityquantity}
  \lean{PhyXMiniProblems.ProblemPhyXMini0238.MassDensityQuantity}
  \uses{def:physics:phyx-mini-0238:phyxminiproblems-problemphyxmini0238-massdensitydimension}
  A nonnegative physical mass density
\end{definition}
\begin{proof}
  This is the declared model definition of Mass Density Quantity.
\end{proof}
\begin{definition}[Mass Flow Rate Quantity]
  \label{def:physics:phyx-mini-0238:phyxminiproblems-problemphyxmini0238-massflowratequantity}
  \lean{PhyXMiniProblems.ProblemPhyXMini0238.MassFlowRateQuantity}
  \uses{def:physics:phyx-mini-0238:phyxminiproblems-problemphyxmini0238-massflowratedimension}
  A nonnegative magnitude of mass discharged per unit time
\end{definition}
\begin{proof}
  This is the declared model definition of Mass Flow Rate Quantity.
\end{proof}
\begin{definition}[Acceleration Quantity]
  \label{def:physics:phyx-mini-0238:phyxminiproblems-problemphyxmini0238-accelerationquantity}
  \lean{PhyXMiniProblems.ProblemPhyXMini0238.AccelerationQuantity}
  \uses{def:physics:phyx-mini-0238:phyxminiproblems-problemphyxmini0238-accelerationdimension}
  A nonnegative physical acceleration
\end{definition}
\begin{proof}
  This is the declared model definition of Acceleration Quantity.
\end{proof}
\begin{definition}[Time Quantity]
  \label{def:physics:phyx-mini-0238:phyxminiproblems-problemphyxmini0238-timequantity}
  \lean{PhyXMiniProblems.ProblemPhyXMini0238.TimeQuantity}
  A nonnegative physical duration
\end{definition}
\begin{proof}
  This is the declared model definition of Time Quantity.
\end{proof}
\begin{definition}[quantity SIReadout]
  \label{def:physics:phyx-mini-0238:phyxminiproblems-problemphyxmini0238-quantitysireadout}
  \lean{PhyXMiniProblems.ProblemPhyXMini0238.quantitySIReadout}
  Read a nonnegative dimensionful quantity in coherent SI units
\end{definition}
\begin{proof}
  This is the declared model definition of quantity SIReadout.
\end{proof}
\begin{definition}[length In Meters]
  \label{def:physics:phyx-mini-0238:phyxminiproblems-problemphyxmini0238-lengthinmeters}
  \lean{PhyXMiniProblems.ProblemPhyXMini0238.lengthInMeters}
  \uses{def:physics:phyx-mini-0238:phyxminiproblems-problemphyxmini0238-lengthquantity, def:physics:phyx-mini-0238:phyxminiproblems-problemphyxmini0238-quantitysireadout}
  Meter readout of a physical length
\end{definition}
\begin{proof}
  This is the declared model definition of length In Meters.
\end{proof}
\begin{definition}[volume In Cubic Meters]
  \label{def:physics:phyx-mini-0238:phyxminiproblems-problemphyxmini0238-volumeincubicmeters}
  \lean{PhyXMiniProblems.ProblemPhyXMini0238.volumeInCubicMeters}
  \uses{def:physics:phyx-mini-0238:phyxminiproblems-problemphyxmini0238-volumequantity, def:physics:phyx-mini-0238:phyxminiproblems-problemphyxmini0238-quantitysireadout}
  Cubic-meter readout of a physical volume
\end{definition}
\begin{proof}
  This is the declared model definition of volume In Cubic Meters.
\end{proof}
\begin{definition}[mass In Kilograms]
  \label{def:physics:phyx-mini-0238:phyxminiproblems-problemphyxmini0238-massinkilograms}
  \lean{PhyXMiniProblems.ProblemPhyXMini0238.massInKilograms}
  \uses{def:physics:phyx-mini-0238:phyxminiproblems-problemphyxmini0238-massquantity, def:physics:phyx-mini-0238:phyxminiproblems-problemphyxmini0238-quantitysireadout}
  Kilogram readout of a physical mass
\end{definition}
\begin{proof}
  This is the declared model definition of mass In Kilograms.
\end{proof}
\begin{definition}[density In Kilograms Per Cubic Meter]
  \label{def:physics:phyx-mini-0238:phyxminiproblems-problemphyxmini0238-densityinkilogramspercubicmeter}
  \lean{PhyXMiniProblems.ProblemPhyXMini0238.densityInKilogramsPerCubicMeter}
  \uses{def:physics:phyx-mini-0238:phyxminiproblems-problemphyxmini0238-massdensityquantity, def:physics:phyx-mini-0238:phyxminiproblems-problemphyxmini0238-quantitysireadout}
  Kilogram-per-cubic-meter readout of a mass density
\end{definition}
\begin{proof}
  This is the declared model definition of density In Kilograms Per Cubic Meter.
\end{proof}
\begin{definition}[mass Flow Rate In Kilograms Per Second]
  \label{def:physics:phyx-mini-0238:phyxminiproblems-problemphyxmini0238-massflowrateinkilogramspersecond}
  \lean{PhyXMiniProblems.ProblemPhyXMini0238.massFlowRateInKilogramsPerSecond}
  \uses{def:physics:phyx-mini-0238:phyxminiproblems-problemphyxmini0238-massflowratequantity, def:physics:phyx-mini-0238:phyxminiproblems-problemphyxmini0238-quantitysireadout}
  Kilogram-per-second readout of the positive outflow-rate magnitude
\end{definition}
\begin{proof}
  This is the declared model definition of mass Flow Rate In Kilograms Per Second.
\end{proof}
\begin{definition}[acceleration In Meters Per Second Squared]
  \label{def:physics:phyx-mini-0238:phyxminiproblems-problemphyxmini0238-accelerationinmeterspersecondsquared}
  \lean{PhyXMiniProblems.ProblemPhyXMini0238.accelerationInMetersPerSecondSquared}
  \uses{def:physics:phyx-mini-0238:phyxminiproblems-problemphyxmini0238-accelerationquantity, def:physics:phyx-mini-0238:phyxminiproblems-problemphyxmini0238-quantitysireadout}
  Meter-per-second-squared readout of a physical acceleration
\end{definition}
\begin{proof}
  This is the declared model definition of acceleration In Meters Per Second Squared.
\end{proof}
\begin{definition}[time In Seconds]
  \label{def:physics:phyx-mini-0238:phyxminiproblems-problemphyxmini0238-timeinseconds}
  \lean{PhyXMiniProblems.ProblemPhyXMini0238.timeInSeconds}
  \uses{def:physics:phyx-mini-0238:phyxminiproblems-problemphyxmini0238-timequantity, def:physics:phyx-mini-0238:phyxminiproblems-problemphyxmini0238-quantitysireadout}
  Second readout of a physical duration
\end{definition}
\begin{proof}
  This is the declared model definition of time In Seconds.
\end{proof}
\begin{definition}[Figure Panel]
  \label{def:physics:phyx-mini-0238:phyxminiproblems-problemphyxmini0238-figurepanel}
  \lean{PhyXMiniProblems.ProblemPhyXMini0238.FigurePanel}
  The two configurations displayed side by side in the source figure
\end{definition}
\begin{proof}
  This is the declared model definition of Figure Panel.
\end{proof}
\begin{definition}[Figure Label]
  \label{def:physics:phyx-mini-0238:phyxminiproblems-problemphyxmini0238-figurelabel}
  \lean{PhyXMiniProblems.ProblemPhyXMini0238.FigureLabel}
  The four labels printed on the source figure
\end{definition}
\begin{proof}
  This is the declared model definition of Figure Label.
\end{proof}
\begin{definition}[Figure Length Role]
  \label{def:physics:phyx-mini-0238:phyxminiproblems-problemphyxmini0238-figurelengthrole}
  \lean{PhyXMiniProblems.ProblemPhyXMini0238.FigureLengthRole}
  Physical meaning of a length marker in the source figure
\end{definition}
\begin{proof}
  This is the declared model definition of Figure Length Role.
\end{proof}
\begin{definition}[length Role]
  \label{def:physics:phyx-mini-0238:phyxminiproblems-problemphyxmini0238-figurelabel-lengthrole}
  \lean{PhyXMiniProblems.ProblemPhyXMini0238.FigureLabel.lengthRole}
  \uses{def:physics:phyx-mini-0238:phyxminiproblems-problemphyxmini0238-figurelabel, def:physics:phyx-mini-0238:phyxminiproblems-problemphyxmini0238-figurelengthrole}
  Meaning assigned to each printed length label
\end{definition}
\begin{proof}
  This is the declared model definition of length Role.
\end{proof}
\begin{definition}[Pendulum Figure]
  \label{def:physics:phyx-mini-0238:phyxminiproblems-problemphyxmini0238-pendulumfigure}
  \lean{PhyXMiniProblems.ProblemPhyXMini0238.PendulumFigure}
  \uses{def:physics:phyx-mini-0238:phyxminiproblems-problemphyxmini0238-figurepanel, def:physics:phyx-mini-0238:phyxminiproblems-problemphyxmini0238-figurelabel}
  Categorical evidence visible in the primary image
\end{definition}
\begin{proof}
  This is the declared model definition of Pendulum Figure.
\end{proof}
\begin{definition}[Outlet Location]
  \label{def:physics:phyx-mini-0238:phyxminiproblems-problemphyxmini0238-outletlocation}
  \lean{PhyXMiniProblems.ProblemPhyXMini0238.OutletLocation}
  The outlet location stated in the prose
\end{definition}
\begin{proof}
  This is the declared model definition of Outlet Location.
\end{proof}
\begin{definition}[Outflow Direction]
  \label{def:physics:phyx-mini-0238:phyxminiproblems-problemphyxmini0238-outflowdirection}
  \lean{PhyXMiniProblems.ProblemPhyXMini0238.OutflowDirection}
  The idealized direction of the expelled liquid relative to the pendulum
\end{definition}
\begin{proof}
  This is the declared model definition of Outflow Direction.
\end{proof}
\begin{definition}[Leaking Cube Pendulum Setup]
  \label{def:physics:phyx-mini-0238:phyxminiproblems-problemphyxmini0238-leakingcubependulumsetup}
  \lean{PhyXMiniProblems.ProblemPhyXMini0238.LeakingCubePendulumSetup}
  \uses{def:physics:phyx-mini-0238:phyxminiproblems-problemphyxmini0238-lengthquantity, def:physics:phyx-mini-0238:phyxminiproblems-problemphyxmini0238-volumequantity, def:physics:phyx-mini-0238:phyxminiproblems-problemphyxmini0238-massquantity, def:physics:phyx-mini-0238:phyxminiproblems-problemphyxmini0238-massdensityquantity, def:physics:phyx-mini-0238:phyxminiproblems-problemphyxmini0238-massflowratequantity, def:physics:phyx-mini-0238:phyxminiproblems-problemphyxmini0238-accelerationquantity, def:physics:phyx-mini-0238:phyxminiproblems-problemphyxmini0238-timequantity, def:physics:phyx-mini-0238:phyxminiproblems-problemphyxmini0238-pendulumfigure, def:physics:phyx-mini-0238:phyxminiproblems-problemphyxmini0238-outletlocation, def:physics:phyx-mini-0238:phyxminiproblems-problemphyxmini0238-outflowdirection}
  All physical quantities in the draining experiment. The trajectory arguments are SI time coordinates in seconds. The supportStringLength runs from the fixed pivot to the top of the cube, whereas both pendulum-length fields run from the pivot to the liquid's center of mass, as specified by the problem statement
\end{definition}
\begin{proof}
  This is the declared model definition of Leaking Cube Pendulum Setup.
\end{proof}
\begin{definition}[Matches Primary Figure]
  \label{def:physics:phyx-mini-0238:phyxminiproblems-problemphyxmini0238-matchesprimaryfigure}
  \lean{PhyXMiniProblems.ProblemPhyXMini0238.MatchesPrimaryFigure}
  \uses{def:physics:phyx-mini-0238:phyxminiproblems-problemphyxmini0238-leakingcubependulumsetup}
  The primary image contains both configurations and all four source labels
\end{definition}
\begin{proof}
  This is the declared model definition of Matches Primary Figure.
\end{proof}
\begin{definition}[Matches Problem Data]
  \label{def:physics:phyx-mini-0238:phyxminiproblems-problemphyxmini0238-matchesproblemdata}
  \lean{PhyXMiniProblems.ProblemPhyXMini0238.MatchesProblemData}
  \uses{def:physics:phyx-mini-0238:phyxminiproblems-problemphyxmini0238-lengthinmeters, def:physics:phyx-mini-0238:phyxminiproblems-problemphyxmini0238-volumeincubicmeters, def:physics:phyx-mini-0238:phyxminiproblems-problemphyxmini0238-massinkilograms, def:physics:phyx-mini-0238:phyxminiproblems-problemphyxmini0238-timeinseconds, def:physics:phyx-mini-0238:phyxminiproblems-problemphyxmini0238-leakingcubependulumsetup}
  Readouts stated in the prose: a cube of volume a³, initially full, with Lᵢ measured to the initial liquid center of mass; a bottom outlet; and complete emptying. Zero container and string masses encode the word "light" in the ideal model. No period value or period approximation occurs here
\end{definition}
\begin{proof}
  This is the declared model definition of Matches Problem Data.
\end{proof}
\begin{definition}[During Emptying]
  \label{def:physics:phyx-mini-0238:phyxminiproblems-problemphyxmini0238-duringemptying}
  \lean{PhyXMiniProblems.ProblemPhyXMini0238.DuringEmptying}
  \uses{def:physics:phyx-mini-0238:phyxminiproblems-problemphyxmini0238-timeinseconds, def:physics:phyx-mini-0238:phyxminiproblems-problemphyxmini0238-leakingcubependulumsetup}
  SI times from release up to, but not including, the massless empty state
\end{definition}
\begin{proof}
  This is the declared model definition of During Emptying.
\end{proof}
\begin{definition}[Has Physical Parameters]
  \label{def:physics:phyx-mini-0238:phyxminiproblems-problemphyxmini0238-hasphysicalparameters}
  \lean{PhyXMiniProblems.ProblemPhyXMini0238.HasPhysicalParameters}
  \uses{def:physics:phyx-mini-0238:phyxminiproblems-problemphyxmini0238-lengthinmeters, def:physics:phyx-mini-0238:phyxminiproblems-problemphyxmini0238-volumeincubicmeters, def:physics:phyx-mini-0238:phyxminiproblems-problemphyxmini0238-densityinkilogramspercubicmeter, def:physics:phyx-mini-0238:phyxminiproblems-problemphyxmini0238-massflowrateinkilogramspersecond, def:physics:phyx-mini-0238:phyxminiproblems-problemphyxmini0238-accelerationinmeterspersecondsquared, def:physics:phyx-mini-0238:phyxminiproblems-problemphyxmini0238-timeinseconds, def:physics:phyx-mini-0238:phyxminiproblems-problemphyxmini0238-leakingcubependulumsetup, def:physics:phyx-mini-0238:phyxminiproblems-problemphyxmini0238-duringemptying}
  Positivity and geometric bounds needed by the mechanical model. The dimensionless ratio bound witnesses the qualitative condition Lᵢ ≫ a
\end{definition}
\begin{proof}
  This is the declared model definition of Has Physical Parameters.
\end{proof}
\begin{definition}[linearized Simple Pendulum Period Seconds]
  \label{def:physics:phyx-mini-0238:phyxminiproblems-problemphyxmini0238-linearizedsimplependulumperiodseconds}
  \lean{PhyXMiniProblems.ProblemPhyXMini0238.linearizedSimplePendulumPeriodSeconds}
  \uses{def:physics:phyx-mini-0238:phyxminiproblems-problemphyxmini0238-lengthquantity, def:physics:phyx-mini-0238:phyxminiproblems-problemphyxmini0238-accelerationquantity, def:physics:phyx-mini-0238:phyxminiproblems-problemphyxmini0238-lengthinmeters, def:physics:phyx-mini-0238:phyxminiproblems-problemphyxmini0238-accelerationinmeterspersecondsquared}
  The period in seconds of the linearized, frozen-length simple-pendulum model. This is an approximating term, not the definition of the measured local period stored in LeakingCubePendulumSetup
\end{definition}
\begin{proof}
  This is the declared model definition of linearized Simple Pendulum Period Seconds.
\end{proof}
\begin{definition}[Obeys Leaking Simple Pendulum Laws]
  \label{def:physics:phyx-mini-0238:phyxminiproblems-problemphyxmini0238-obeysleakingsimplependulumlaws}
  \lean{PhyXMiniProblems.ProblemPhyXMini0238.ObeysLeakingSimplePendulumLaws}
  \uses{def:physics:phyx-mini-0238:phyxminiproblems-problemphyxmini0238-lengthinmeters, def:physics:phyx-mini-0238:phyxminiproblems-problemphyxmini0238-volumeincubicmeters, def:physics:phyx-mini-0238:phyxminiproblems-problemphyxmini0238-massinkilograms, def:physics:phyx-mini-0238:phyxminiproblems-problemphyxmini0238-densityinkilogramspercubicmeter, def:physics:phyx-mini-0238:phyxminiproblems-problemphyxmini0238-massflowrateinkilogramspersecond, def:physics:phyx-mini-0238:phyxminiproblems-problemphyxmini0238-timeinseconds, def:physics:phyx-mini-0238:phyxminiproblems-problemphyxmini0238-leakingcubependulumsetup, def:physics:phyx-mini-0238:phyxminiproblems-problemphyxmini0238-duringemptying, def:physics:phyx-mini-0238:phyxminiproblems-problemphyxmini0238-linearizedsimplependulumperiodseconds}
  The liquid is uniform and fills a square horizontal cross-section. Its volume is a² h, its mass is density times volume, and its positive constant outflow-rate magnitude gives M(t) = M(0) - q t up to the emptying time. The center-of-mass equations use that a uniform liquid column has its center at half its fill height.
\end{definition}
\begin{proof}
  This is the declared model definition of Obeys Leaking Simple Pendulum Laws.
\end{proof}
\begin{definition}[Answer Choice]
  \label{def:physics:phyx-mini-0238:phyxminiproblems-problemphyxmini0238-answerchoice}
  \lean{PhyXMiniProblems.ProblemPhyXMini0238.AnswerChoice}
  Labels of the four answer choices in the source
\end{definition}
\begin{proof}
  This is the declared model definition of Answer Choice.
\end{proof}
\begin{definition}[displayed Scalar]
  \label{def:physics:phyx-mini-0238:phyxminiproblems-problemphyxmini0238-answerchoice-displayedscalar}
  \lean{PhyXMiniProblems.ProblemPhyXMini0238.AnswerChoice.displayedScalar}
  \uses{def:physics:phyx-mini-0238:phyxminiproblems-problemphyxmini0238-answerchoice}
  The scalar expressions printed beside the four choices after substituting SI readouts for Lᵢ and g. This records the source display; in particular, some distractors do not have the physical dimension of time
\end{definition}
\begin{proof}
  This is the declared model definition of displayed Scalar.
\end{proof}
\begin{definition}[recorded Answer Choice]
  \label{def:physics:phyx-mini-0238:phyxminiproblems-problemphyxmini0238-recordedanswerchoice}
  \lean{PhyXMiniProblems.ProblemPhyXMini0238.recordedAnswerChoice}
  \uses{def:physics:phyx-mini-0238:phyxminiproblems-problemphyxmini0238-answerchoice}
  The dataset records answer D
\end{definition}
\begin{proof}
  This is the declared model definition of recorded Answer Choice.
\end{proof}
\begin{definition}[small Angle Period Error Budget Seconds]
  \label{def:physics:phyx-mini-0238:phyxminiproblems-problemphyxmini0238-smallangleperioderrorbudgetseconds}
  \lean{PhyXMiniProblems.ProblemPhyXMini0238.smallAnglePeriodErrorBudgetSeconds}
  \uses{def:physics:phyx-mini-0238:phyxminiproblems-problemphyxmini0238-leakingcubependulumsetup, def:physics:phyx-mini-0238:phyxminiproblems-problemphyxmini0238-linearizedsimplependulumperiodseconds}
  The explicit small-angle contribution to the local-period error budget. It vanishes quadratically with the dimensionless oscillation amplitude
\end{definition}
\begin{proof}
  This is the declared model definition of small Angle Period Error Budget Seconds.
\end{proof}
\begin{definition}[long String Period Error Budget Seconds]
  \label{def:physics:phyx-mini-0238:phyxminiproblems-problemphyxmini0238-longstringperioderrorbudgetseconds}
  \lean{PhyXMiniProblems.ProblemPhyXMini0238.longStringPeriodErrorBudgetSeconds}
  \uses{def:physics:phyx-mini-0238:phyxminiproblems-problemphyxmini0238-lengthinmeters, def:physics:phyx-mini-0238:phyxminiproblems-problemphyxmini0238-leakingcubependulumsetup, def:physics:phyx-mini-0238:phyxminiproblems-problemphyxmini0238-linearizedsimplependulumperiodseconds}
  The long-string error budget obtained from the center-of-mass displacement 0 ≤ L(t) - Lᵢ ≤ a/2 and concavity of the square root. Relative to the choice-D leading term it is first order in the dimensionless ratio a/Lᵢ
\end{definition}
\begin{proof}
  This is the declared model definition of long String Period Error Budget Seconds.
\end{proof}
% --- Archon declaration topology end ---
% --- Archon physics formalization source end ---
